
\section*{\underline{EXPOS\'E XIII}\\[3pt] \underline{\underline{LES TH\'EOR\`EMES DE FINITUDE POUR LE FONCTEUR DE PICARD}}}
\begin{center}
\underline{par S. KLEIMAN}
\end{center}

\subsection*{1. Les (b)-faisceaux}
Soient $k$ un corps alg\'ebriquement clos et $X$ un $k$-sch\'ema projectif, muni d'un module inversible ample ${O}_X(1)$.

\textbf{D\'efinition 1.1 (Mumford).} \emph{Soit $m\in\mathbb Z$. Un faisceau coh\'erent $F$ sur $X$ est dit \emph{$m$-r\'egulier} si, en chaque point du support de $F$, ${O}_X(1)$ est engendr\'e par ses sections globales et si, pour tout $q\ge1$, on a}
\[
 H^q(F(m-q))=0.
\]

\textbf{Lemme 1.2.} \emph{Soit \(0\longrightarrow F(-1)\longrightarrow F\longrightarrow G\longrightarrow 0\) une suite exacte de faisceaux coh\'erents sur $X$}\SGAAuthorFootnote{*}{Il est sous-entendu, ici comme dans la suite dans des situations analogues, que \(F(-1)\to F\) est d\'efini par tensorisation par une section donn\'ee de \(\underline{O}_X(1)\).}\emph{. Si $F$ est $m$-r\'egulier, alors $G$ l'est.}

En effet, on a la suite exacte correspondante :
\[
H^q(F(m-q))\longrightarrow H^q(G(m-q))\longrightarrow H^{q+1}(F(m-q-1)).
\]

\textbf{Proposition 1.3.} \emph{Si un faisceau $F$ sur $X$ est $m$-r\'egulier, alors, pour tout $n\ge m$, on a}~:
\begin{enumerate}[label=(\roman*)]
\item \emph{$F$ est $n$-r\'egulier.}
\item \emph{\(H^0(F(n))\otimes H^0({O}_X(1))\longrightarrow H^0(F(n+1))\) est surjectif.}
\item \emph{$F(n)$ est engendr\'e par \(H^0(F(n))\).}
\end{enumerate}

